% Baccalauréat malgache, série D, 2009 : énoncé.
\TCDBacAttribution{D}{2009}
\providecommand{\Courbe}{(\mathscr{C})}
\providecommand{\D}{(\mathscr{D})}
\providecommand{\Plan}{(\mathscr{P})}
\providecommand{\Ouv}{(O ; \vec{u} , \vec{v})}
%\providecommand{\Oij}{(O ; \vec{\imath} , \vec{\jmath})}
\providecommand{\cm}{~\mathrm{cm}}
\providecommand{\dx}{~\mathrm{d}x}
\providecommand{\Arg}[1]{\mathrm{Arg}\left( #1 \right)}
%%%%%%%%%%%%%%%%%%%%%%%%%%%%%%%%%%%%%%%%%%%%%%%%%%%%%%%%%%%%%%%%%%%%%%%%%%%%%%%%%%%%%%%%%%%%%%%
%%%%%%%%%%%%%%%%%%%                   %%%%%%%%%%%%%%%%%%%                   %%%%%%%%%%%%%%%%%%%
%%%%%%%%%%%%%%%%%%%%%%%%%%%%%%%%%%%%%%%%%%%%%%%%%%%%%%%%%%%%%%%%%%%%%%%%%%%%%%%%%%%%%%%%%%%%%%%
\noindent \underline{\textsc{Exercice I}} (5 points) \newline
      Le plan complexe $\Plan$ est muni d'un rep\`ere orthonorm\'e $\Ouv$ d'unit\'e $1\cm$.
      \begin{enumerate}
            \item $P$ est le polyn\^ome de la variable complexe $z$ d\'efini par : 
                  $P(z) = z^3 - (6+2i) z^2 + (10+8i) z - 4 - 8i$.
                  \begin{enumerate}
                        \item Calculer $P(2)$.
                        \item R\'esoudre dans $\mathbb{C}$, l'\'equation $P(z) = 0$.
                  \end{enumerate}
            \item On donne les points $A$, $B$, $C$ et $D$ d'affixes respectives $a = 1+i$ ; $b = 2$ ; $c = 3+i$ et $d$.
                  \begin{enumerate}
                        \item Calculer $d$ pour que $ABCD$ soit un carr\'e.
                        \item On consid\`ere la similitude plane directe $S$ d\'efinie par son expression complexe :
                              $z' = (1+i) z + 4i$.
                              \begin{itemize}
                                    \item Donner les \'el\'ements caract\'eristiques de $S$.
                                    \item D\'eterminer l'expression analytiques de $S$.
                              \end{itemize}
                  \end{enumerate}
            \item Construire dans le m\^eme rep\`ere $ABCD$ et son image $A'B'C'D'$ par $S$.
      \end{enumerate}
%%%%%%%%%%%%%%%%%%%%%%%%%%%%%%%%%%%%%%%%%%%%%%%%%%%%%%%%%%%%%%%%%%%%%%%%%%%%%%%%%%%%%%%%%%%%%%%
%%%%%%%%%%%%%%%%%%%                   %%%%%%%%%%%%%%%%%%%                   %%%%%%%%%%%%%%%%%%%
%%%%%%%%%%%%%%%%%%%%%%%%%%%%%%%%%%%%%%%%%%%%%%%%%%%%%%%%%%%%%%%%%%%%%%%%%%%%%%%%%%%%%%%%%%%%%%%
\noindent \underline{\textsc{Exercice II}} (5 points) 
      \begin{enumerate}
            \item Les faces d'un d\'e cubique $D_1$ truqu\'e sont num\'erot\'ees : 1, 2, 2, 3, 3, 3.
                  On lance une fois ce d\'e.
                  Chaque face a la m\^eme probabilit\'e d'apparition.
                  On note $P_i$ la probabilit\'e d'apparition de la face portant le num\'ero $i$.
                  Calculer $P_1$, $P_2$ et $P_3$.
            \item Les faces d'un deuxi\`eme d\'e cubique normal $D_2$ sont num\'erot\'ees de 1 \`a 6.
                  On lance en m\^eme temps les deux d\'es $D_1$ et $D_2$.
                  Chaque face a toujours la m\^eme probabilit\'e d'apparition.
                  \begin{enumerate}
                        \item Calculer la probabilit\'e de l'\'ev\`enement : 
                              $A$ : << Les deux d\'es affichent le m\^eme num\'ero >>.
                        \item On d\'esigne par $X$ la variable al\'eatoire d\'efinie par la somme des num\'eros affich\'es
                              par les deux d\'es.
                              \begin{itemize}
                                    \item Donner la loi de probabilit\'e de $X$.
                                    \item Calculer l'esp\'erance math\'ematique $E(X)$.
                              \end{itemize}
                        \item On lance trois fois de suite et d'une fa\c{c}on ind\'ependante le d\'e $D_1$.
                              On note $Y$ la variable al\'eatoire \'egale au 
                              nombre d'apparitions de la face portant le num\'ero 2 lors de ces trois lancers.
                              \begin{enumerate}
                                    \item Donner la loi de probabilit\'e de $Y$.
                                    \item Calculer la variance $V(Y)$.
                              \end{enumerate}
                  \end{enumerate}
            \begin{center}
                  NB : Tous les r\'esultats seront exprim\'es sous forme de fraction irr\'eductible.
            \end{center}
      \end{enumerate}
%%%%%%%%%%%%%%%%%%%%%%%%%%%%%%%%%%%%%%%%%%%%%%%%%%%%%%%%%%%%%%%%%%%%%%%%%%%%%%%%%%%%%%%%%%%%%%%
%%%%%%%%%%%%%%%%%%%                   %%%%%%%%%%%%%%%%%%%                   %%%%%%%%%%%%%%%%%%%
%%%%%%%%%%%%%%%%%%%%%%%%%%%%%%%%%%%%%%%%%%%%%%%%%%%%%%%%%%%%%%%%%%%%%%%%%%%%%%%%%%%%%%%%%%%%%%%
\noindent \underline{\textsc{Probleme}} \newline
      On consid\`ere la fonction num\'erique $f$ de la variable r\'eelle $x$ d\'efinie par : 
            \[
                  \begin{multiform}
                        f(0) = 0 \\
                        f(x) = x (1-\ln x)^2 \text{ si } x > 0 
                  \end{multiform}
            \]
            o\`u $\ln$ d\'esigne la fonction logarithme n\'eperien.
      On appelle $\Courbe$ sa courbe repr\'esentative dans un rep\`ere orthonorm\'e $R = \text{\Oij}$ d'unit\'e $1\cm$.
      \begin{enumerate}
            \item Justifier que l'ensemble de d\'efinition de $f$ est $[0;+\infty[$.
            \item \begin{enumerate}
                        \item Prouver que $f$ est continue en $x_0 = 0$.
                        \item Etudier la d\'erivabilit\'e de $f$ en $x_0 = 0$.
                  \end{enumerate}
            \item \begin{enumerate}
                        \item D\'emontrer que pour tout $x>0$ ; $f'(x) = (\ln x - 1) (\ln x + 1)$ o\`u $f'$ d\'esigne la fonction
                              d\'eriv\'ee premi\`ere de $f$.
                        \item Apr\`es avoir pr\'eciser le sens de variations de $f$, dresser son tableau de variations.
                  \end{enumerate}
            \item Tracer la courbe $\Courbe$.
            \item Soit $\alpha$ un r\'eel appartenant \`a l'intervalle $]0;e[$.
                  \begin{enumerate}
                        \item Calculer, \`a l'aide de deux int\'egrations par parties l'int\'egrale 
                              \[
                                    I = \int_{\alpha}^{e} f(x) \dx \;.
                              \]
                        \item Soit $A(\alpha)$ l'aire du domaine plan d\'elimit\'e par l'axe $(O;\ivec)$,
                              les droites d'\'equations $x = \alpha$ et $x = e$ et la courbe $\Courbe$.
                              Calculer en $\cm^2$ $A(\alpha)$ et $\displaystyle{\lim_{\alpha \rightarrow 0^+} A(\alpha)}$.
                  \end{enumerate}
            \item Soit $g$ la restriction de $f$ \`a l'intervalle $[e;+\infty[$.
                  \begin{enumerate}
                        \item D\'emontrer que $g$ admet une application r\'eciproque not\'ee $g^{-1}$ dont on pr\'ecisera
                              l'ensemble de d\'efinition.
                        \item Calculer $g(\expo{2})$ et $\left( g^{-1} \right)' (\expo{2})$.
                  \end{enumerate}
            \item Tracer dans le m\^eme rep\`ere que $\Courbe$ la courbe repr\'esentative $(\Gamma)$ de $g^{-1}$.
      \end{enumerate}
            %
            On donne $e \simeq 2,7$ ; $\expo{-1} \simeq 0,36$ ; $\expo{2} \simeq 7,4$.











